\documentclass[11pt]{article}
\usepackage[a4paper,margin=18mm]{geometry}
\usepackage[no-math]{fontspec}
\setmainfont{Latin Modern Roman}
\usepackage{amsmath,amssymb,amsthm,xfrac,xspace,hyperref,graphicx,comment}
\excludecomment{DefTralics}

\providecommand{\bibliofont}{\small}
\newcommand{\ie}{i.e.,\xspace}
\newcommand{\eg}{e.g.,\xspace}
\newcommand{\etc}{etc.\xspace}
\newcommand{\cf}{cf.\xspace}
\newcommand{\resp}{resp.\xspace}
\newcommand{\smaller}{\small}
\newcommand{\Subsection}{\subsection}
\newcommand{\Subsubsection}{\subsubsection}
\newcommand{\skpt}{\smallskip}
\emergencystretch=3em
\numberwithin{equation}{section}\input{author-preamble.tex}
\makeatletter
\expandafter\def\csname jeplabeltype@thm:main-inter-L2\endcsname{theorem}
\expandafter\def\csname jeplabeltype@prop:co-eigenfunctions\endcsname{proposition}
\expandafter\def\csname jeplabeltype@thm:spectral-representation\endcsname{theorem}
\expandafter\def\csname jeplabeltype@thm:convergence-equilibrium\endcsname{theorem}
\expandafter\def\csname jeplabeltype@thm:contractivity\endcsname{theorem}
\expandafter\def\csname jeplabeltype@cor:subordinated\endcsname{corollary}
\expandafter\def\csname jeplabeltype@lem:rewriting-stuff\endcsname{lemma}
\expandafter\def\csname jeplabeltype@lem:phi-Psi1-claims\endcsname{lemma}
\expandafter\def\csname jeplabeltype@lem:Markov-kernel\endcsname{lemma}
\expandafter\def\csname jeplabeltype@lem:fact\endcsname{lemma}
\expandafter\def\csname jeplabeltype@rem:facto\endcsname{remark}
\expandafter\def\csname jeplabeltype@prop:inter_P\endcsname{proposition}
\expandafter\def\csname jeplabeltype@lem:inter\endcsname{lemma}
\expandafter\def\csname jeplabeltype@lem:op-com\endcsname{lemma}
\expandafter\def\csname jeplabeltype@lem:feller\endcsname{lemma}
\expandafter\def\csname jeplabeltype@prop:Convolution\endcsname{proposition}
\expandafter\def\csname jeplabeltype@lem:Mellin-bpsi\endcsname{lemma}
\expandafter\def\csname jeplabeltype@prop:inter_P-2\endcsname{proposition}
\expandafter\def\csname jeplabeltype@prop:intertwining\endcsname{proposition}
\expandafter\def\csname jeplabeltype@lem:Lambda-L2-bounded\endcsname{lemma}
\expandafter\def\csname jeplabeltype@prop:resolvent-intertwining\endcsname{proposition}
\expandafter\def\csname jeplabeltype@lem:seq-cn\endcsname{lemma}
\expandafter\def\csname jeplabeltype@prop:riesz\endcsname{proposition}
\expandafter\def\csname jeplabeltype@rem:bounds\endcsname{remark}
\expandafter\def\csname jeplabeltype@prop:quasi-similar\endcsname{proposition}
\expandafter\def\csname jeplabeltype@prop:cmir\endcsname{proposition}
\expandafter\def\csname jeplabeltype@prop:small-perturbation\endcsname{proposition}
\AtBeginDocument{\let\jeporiginalLabel\label
\renewcommand{\label}[1]{\ifcsname jeplabeltype@#1\endcsname
\edef\jeplabeltype{\csname jeplabeltype@#1\endcsname}\expandafter\jeporiginalLabel\expandafter[\jeplabeltype]{#1}\else\jeporiginalLabel{#1}\fi}}
\makeatother
\begin{document}
\begin{center}{\Large Stable item 2045}\end{center}
\paragraph{Source and attribution.}
Patrick Cheridito, Pierre Patie, Anna Srapionyan, Aditya Vaidyanathan, \emph{On non-local ergodic Jacobi semigroups: spectral theory, convergence-to-equilibrium and contractivity}, Journal de l’École polytechnique — Mathématiques 8 (2021), publisher version, DOI 10.5802/jep.148. \href{https://jep.centre-mersenne.org/articles/10.5802/jep.148/}{Publisher article}. Authors' text: \href{https://creativecommons.org/licenses/by/4.0/}{CC BY 4.0}.
\paragraph{Editorial problem boundary.}
Prove only Theorem 2.1 below for a genuinely nonzero downward-jump kernel satisfying every printed Assumption A condition, including the finite positive Radon jump measure and exact drift/parameter inequalities. The moment-determinate invariant law, continuous positive interior density and ergodic Markov-semigroup generation with polynomial core are selected. Full new intertwining, maximum-principle, polynomial resolvent and invariant-measure uniqueness core is retained. The classical zero-kernel branch, spectral expansion and entropy-rate assertions are context only.
\paragraph{Difficulty (editorial): H2.}
The classical Jacobi semigroup and Bernstein-function prerequisites do not themselves construct the nonlocal generator with the stated positive invariant density. The author proves the necessary moment factorisation and full support, builds an explicit polynomial intertwining and maximum-principle/resolvent argument, and transfers the invariant measure without assuming the resulting Markov process exists. Those constructions then establish the polynomial core, invariant-law uniqueness and ergodicity in the genuinely nonzero jump-kernel case.
\paragraph{Editorial transcription note.}
Original author language, mathematics and ordering are preserved. Publisher front matter is replaced by this independent article wrapper. Bibliography is recovered from matching cached publisher HTML, including its TeX formulas. No new proof or mathematical repair is supplied. Surrounding original results are retained for conservative context and are not extra selected corpus claims.
\clearpage
\input{author-body.tex}
\begin{thebibliography}{99}
\bibitem{BGK} F. Achleitner, A. Arnold \& E. A. Carlen - “On multi-dimensional hypocoercive BGK models” , Kinet. and Relat. Mod. 11 (2018) no. 4, p. 953-1009
\bibitem{logSobolev-book} C. Ané, S. Blachère, D. Chafaï, P. Fougères, I. Gentil, F. Malrieu, C. Roberto \& G. Scheffer - Sur les inégalités de Sobolev logarithmiques , Panoramas \& Synthèses , vol. 10 , Société Mathématique de France, Paris, 2000
\bibitem{Bakry-Jacobi} D. Bakry - “Remarques sur les semigroupes de Jacobi” , in Hommage à P. A. Meyer et J. Neveu , Astérisque , vol. 236 , Société Mathématique de France, Paris, 1996, p. 23-39
\bibitem{Bakry-relativistic} D. Bakry - “Étude des transformations de Riesz dans les variétés riemanniennes à courbure de Ricci minorée” , in Séminaire de Probabilités, XXI , Lect. Notes in Math. , vol. 1247 , Springer, Berlin, 1987, p. 137-172
\bibitem{Bakry_Book} D. Bakry, I. Gentil \& M. Ledoux - Analysis and geometry of Markov diffusion operators , Grundlehren Math. Wiss. , vol. 348 , Springer, Cham, 2014
\bibitem{Baudoin} F. Baudoin - “Bakry-Émery meet Villani” , J. Funct. Anal. 273 (2017) no. 7, p. 2275-2291
\bibitem{Berg} C. Berg \& A. J. Durán - “A transformation from Hausdorff to Stieltjes moment sequences” , Ark. Mat. 42 (2004) no. 2, p. 239-257
\bibitem{Bertoin-s} J. Bertoin - “Subordinators: examples and applications” , in Lectures on probability theory and statistics (Saint-Flour, 1997) , Lect. Notes in Math. , vol. 1717 , Springer, Berlin, 1999, p. 1-91
\bibitem{Potential2009} K. Bogdan, T. Byczkowski, T. Kulczycki, M. Ryznar, R. Song \& Z. Vondraček - Potential analysis of stable processes and its extensions , Lect. Notes in Math. , vol. 1980 , Springer-Verlag, Berlin, 2009
\bibitem{thoma-cone} A. Borodin \& G. Olshanski - “Markov dynamics on the Thoma cone: a model of time-dependent determinantal processes with infinitely many particles” , Electron. J. Probab. 18 (2013), article ID 75, 43 pages
\bibitem{Borodin2002} A. N. Borodin \& P. Salminen - Handbook of Brownian motion—facts and formulae , Probability and its Applications , Birkhäuser Verlag, Basel, 2002
\bibitem{Levy-Matters-III} B. Böttcher, R. Schilling \& J. Wang - Lévy. III , Lect. Notes in Math. , vol. 2099 , Springer, Cham, 2013
\bibitem{barnes_integrals} B. L. J. Braaksma - “Asymptotic expansions and analytic continuations for a class of Barnes-integrals” , Compositio Math. 15 (1964), p. 239-341 (1964)
\bibitem{chafai} D. Chafaï - “Entropies, convexity, and functional inequalities: on $\Phi $-entropies and $\Phi $-Sobolev inequalities” , J. Math. Kyoto Univ. 44 (2004) no. 2, p. 325-363
\bibitem{Chazal2014} M. Chazal, A. Kyprianou \& P. Patie - “A transformation for Lévy processes with one-sided jumps with applications” , 2010
\bibitem{Christensen2003a} O. Christensen - An introduction to frames and Riesz bases , Applied and Numerical Harmonic Analysis , Birkhäuser/Springer, Cham, 2016
\bibitem{cuchiero2018} C. Cuchiero, M. Larsson \& S. Svaluto-Ferro - “Polynomial jump-diffusions on the unit simplex” , Ann. Appl. Probab. 28 (2018) no. 4, p. 2451-2500
\bibitem{da-prato:2006} G. Da Prato - An introduction to infinite-dimensional analysis , Universitext , Springer-Verlag, Berlin, 2006
\bibitem{delbaen2002interest} F. Delbaen \& H. Shirakawa - “An Interest Rate Model with Upper and Lower Bounds” , Asia-Pacific Financial Markets 9 (2002) no. 3, p. 191-209
\bibitem{dem} N. Demni \& M. Zani - “Large deviations for statistics of the Jacobi process” , Stochastic Process. Appl. 119 (2009) no. 2, p. 518-533
\bibitem{Mouhot} J. Dolbeault, C. Mouhot \& C. Schmeiser - “Hypocoercivity for linear kinetic equations conserving mass” , Trans. Amer. Math. Soc. 367 (2015) no. 6, p. 3807-3828
\bibitem{Driver03} B. K. Driver - “Analysis tools with applications” (2003), http://www.math.ucsd.edu/\textasciitilde{}bdriver/231-02-03/Lecture\_Notes/PDE-Anal-Book/analpde1.pdf
\bibitem{dynkin:1965} E. B. Dynkin - Markov processes. Vols. I, II , Grundlehren Math. Wiss. , vol. 121; 122 , Academic Press Inc., New York; Springer-Verlag, Berlin-Göttingen-Heidelberg, 1965
\bibitem{EngelNagel} K.-J. Engel \& R. Nagel - One-parameter semigroups for linear evolution equations , Graduate Texts in Math. , vol. 194 , Springer-Verlag, New York, 2000
\bibitem{EthierKurtz} S. N. Ethier \& T. G. Kurtz - Markov processes , Wiley Series in Proba. and Math. Stat. , John Wiley \& Sons, Inc., New York, 1986
\bibitem{fontenas:1998} É. Fontenas - “Sur les minorations des constantes de Sobolev et de Sobolev logarithmiques pour les opérateurs de Jacobi et de Laguerre” , in Séminaire de Probabilités, XXXII , Lect. Notes in Math. , vol. 1686 , Springer, Berlin, 1998, p. 14-29
\bibitem{gourieroux2006} C. Gourieroux \& J. Jasiak - “Multivariate Jacobi process with application to smooth transitions” , J. Econometrics 131 (2006) no. 1-2, p. 475-505
\bibitem{griffiths2018} R. C. Griffiths, P. A. Jenkins \& D. Spanó - “Wright-Fisher diffusion bridges” , Theoret. Population Biol. 122 (2018), p. 67-77
\bibitem{griffiths:2010} R. C. Griffiths \& D. Spanó - “Diffusion processes and coalescent trees” , in Probability and mathematical genetics , London Math. Soc. Lecture Note Ser. , vol. 378 , Cambridge Univ. Press, Cambridge, 2010, p. 358-379
\bibitem{Gross} L. Gross - “Logarithmic Sobolev inequalities” , Amer. J. Math. 97 (1975) no. 4, p. 1061-1083
\bibitem{huillet2007wright} T. Huillet - “On Wright-Fisher diffusion and its relatives” , J. Stat. Mech. Theory Exp. (2007) no. 11, article ID P11006, 30 pages
\bibitem{Ismail} M. E. H. Ismail - Classical and quantum orthogonal polynomials in one variable , Encyclopedia of Math. and its Appl. , vol. 98 , Cambridge University Press, Cambridge, 2005
\bibitem{Kyprianou2014} A. E. Kyprianou - Fluctuations of Lévy processes with applications , Universitext , Springer, Heidelberg, 2014
\bibitem{Gateway} L. Miclo \& P. Patie - “On a gateway between continuous and discrete Bessel and Laguerre processes” , Ann. H. Lebesgue 2 (2019), p. 59-98
\bibitem{Patie-Miclo} L. Miclo \& P. Patie - “On interweaving relations” , J. Funct. Anal. 280 (2021) no. 3, p. 108816, 53
\bibitem{Mischler2016} S. Mischler \& C. Mouhot - “Exponential stability of slowly decaying solutions to the kinetic-Fokker-Planck equation” , Arch. Rational Mech. Anal. 221 (2016) no. 2, p. 677-723
\bibitem{Misra-Lavoine} O. P. Misra \& J. L. Lavoine - Transform analysis of generalized functions , North-Holland Math. Studies , vol. 119 , North-Holland Publishing Co., Amsterdam, 1986
\bibitem{pal2013} S. Pal - “Wright-Fisher diffusion with negative mutation rates” , Ann. Probab. 41 (2013) no. 2, p. 503-526
\bibitem{Patie-11-Factorization_e} P. Patie - “A refined factorization of the exponential law” , Bernoulli 17 (2011) no. 2, p. 814-826
\bibitem{Patie-Savov-Bern} P. Patie \& M. Savov - “Bernstein-gamma functions and exponential functionals of Lévy processes” , Electron. J. Probab. 23 (2018), article ID 75, 101 pages
\bibitem{Patie-Savov-GeL} P. Patie \& M. Savov - Spectral expansion of non-self-adjoint generalized Laguerre semigroups , Mem. Amer. Math. Soc. , American Mathematical Society, Providence, RI, 2019, to appear
\bibitem{Intertwining-Krein} P. Patie, M. Savov \& Y. Zhao - “Intertwining, excursion theory and Krein theory of strings for non-self-adjoint Markov semigroups” , Ann. Probab. 47 (2019) no. 5, p. 3231-3277
\bibitem{PV-Hypo} P. Patie \& A. Vaidyanathan - “A spectral theoretical approach for hypocoercivity applied to some degenerate hypoelliptic, and non-local operators” , Kinet. and Relat. Mod. 13 (2020) no. 3, p. 479-506
\bibitem{pazy:1983} A. Pazy - Semigroups of linear operators and applications to partial differential equations , Applied Math. Sciences , vol. 44 , Springer-Verlag, New York, 1983
\bibitem{pearson:1992} J. M. Pearson - “Best constants in Sobolev inequalities for ultraspherical polynomials” , Arch. Rational Mech. Anal. 116 (1992) no. 4, p. 361-374
\bibitem{rogers:1981} L. C. G. Rogers \& J. W. Pitman - “Markov functions” , Ann. Probab. 9 (1981) no. 4, p. 573-582
\bibitem{Rudin1973} W. Rudin - Functional analysis , McGraw-Hill Book Co., New York-Düsseldorf-Johannesburg, 1973
\bibitem{Sato1999} K.-i. Sato - Lévy processes and infinitely divisible distributions , Cambridge Studies in Advanced Math. , vol. 68 , Cambridge University Press, Cambridge, 1999
\bibitem{SchillingSongVondracek10} R. L. Schilling, R. Song \& Z. Vondraček - Bernstein functions , De Gruyter Studies in Math. , vol. 37 , Walter de Gruyter \& Co., Berlin, 2012
\bibitem{orthogonal_pols} G. Szegő - Orthogonal polynomials , American Math. Soc., Colloquium Publ. , vol. XXIII , American Mathematical Society, Providence, RI, 1975
\bibitem{Villani-09} C. Villani - Hypocoercivity , Mem. Amer. Math. Soc. , vol. 202, no. 950 , American Mathematical Society, Providence, RI, 2009
\end{thebibliography}
\end{document}
